% !TEX TS-program = xelatex
% !BIB program = bibtex
% !TeX spellcheck = ru_RU
% !TEX root = talk.tex

\documentclass
  [ russian
  , aspectratio=43
  ] {beamer}


\input{preamble.tex}
\makeatletter

% ------------------------------------------

\usepackage{xcolor}
\setbeamercolor{item}{fg=red}

\usepackage{listings}
\usepackage{listingsutf8}

\usepackage[dvipsnames]{xcolor}

\lstdefinelanguage{ML}{
  keywords={type, match, with, let, rec, in, fun, function, of, if, then, else},
  keywordstyle=\color{Maroon}\bfseries,
  identifierstyle=\color{black},
  comment=[(*][*)],
  commentstyle=\color{gray},
  stringstyle=\color{red},
}

\lstset{
    basicstyle=\ttfamily\small
}

\setmonofont{DejaVu Sans Mono}

% ------------------------------------------

\input{pretitle.tex}
\input{title.tex}

\newcommand{\academicGroup}{\my@title@group@ru}
\newcommand{\advisorChair}{\my@title@chair@ru}
\title[Добавление ADT в rukaml]{\my@title@title@ru}
\author[\my@title@author@ru]{\my@title@author@ru, группа \academicGroup}
\institute[СПбГУ]{}
\date[29 декабря 2025 г.]{}
\newcommand{\supervisor}{\my@title@supervisor@ru}
\newcommand{\supervisorPosition}{\my@title@supervisorPosition@ru}
\newcommand{\consultant}{\my@title@consultant@ru}
\newcommand{\consultantPosition}{\my@title@consultantPosition@ru}
\newcommand{\reviewer}{\my@title@reviewer@ru}
\newcommand{\reviewerPosition}{\my@title@reviewerPosition@ru}
\newcommand{\defenseYear}{\my@title@year@ru}

\makeatother
\begin{document}
{
\setbeamertemplate{footline}{}
\begin{frame}
    \includegraphics[width=1.4cm]{figures/SPbGU_Logo.png}
    \vspace{-35pt}
    \hspace{-10pt}
    \begin{center}
        \begin{tabular}{c}
            \scriptsize{Санкт-Петербургский государственный университет} \\
            \scriptsize{\advisorChair}
        \end{tabular}
        \titlepage
    \end{center}

    \btVFill

    {\scriptsize
        \textbf{Научный руководитель:}  \supervisorPosition~\supervisor \\
    }
    \makeatother
    \begin{center}
        \vspace{5pt}
        \scriptsize{Санкт-Петербург\\ \defenseYear}
    \end{center}
\end{frame}
}


\begin{frame}
    \frametitle{Постановка задачи}
    \textbf{Целью} работы является реализация поддержки алгебраических типов данных и паттерн-матчинга выражений в rukaml
    \vspace{1em}

    \textbf{Задачи}:
    \begin{itemize}
        \item Добавить в парсер обработку новых конструкций
        \item Добавить поддержку списков как частного случая ADT в качестве built-in типа
        \item Изменить окружение тайпчекера таким образом, чтобы можно было выводить ADT
        \item Написать порождение кода для значений алгебраических типов
    \end{itemize}
\end{frame}


\begin{frame}[fragile]
    \frametitle{Новые конструкции}
    Объявления алгебраических типов данных и их конструкторы

    {
    \begin{lstlisting}[language=ML]
    type ('a1, ..., 'aN) type_name =
      | Constructor_name1 of ...
      | Constructor_name2 of ...
      | ...
      | Constructor_nameM of ...
    \end{lstlisting}

    Сопоставление с образцом
    \begin{lstlisting}[language=ML]
    match expression with
    | pattern1 -> expression1
    | pattern2 -> expression2
    | ...
    | patternN -> expressionN
    \end{lstlisting}

    }
\end{frame}


\begin{frame}[fragile]
    \frametitle{Списки}

    Синтаксический сахар для конструкторов и их обработка в pretty-printer'ах
    \begin{lstlisting}[language=ML]
    x :: [ a; b; c ]

    let rec take n ls =
      match ls with
      | [] -> []
      | x :: xs ->
        if n < 1 then [] else x :: take (n - 1) xs
    \end{lstlisting}

    В начальном окружении тайпчекера описываются встроенные типы и их конструкторы, в том числе список
    \begin{lstlisting}[language=ML]

    let typ_list : type_declaration = ...
    let constr_nil : constructor_info = ...
    let constr_cons : constructor_info = ...
    \end{lstlisting}
\end{frame}


\begin{frame}
    \frametitle{Вывод типов}
    Нужно научить тайпчекер
    \begin{itemize}
        \item Проверять, что используемые типы уже были объявлены
        \item Проверять, что ожидаемый тип аргумента конструктора унифицируется с передаваемым
        \item По имени конструктора находить тип, в рамках которого он был объявлен и унифицировать его полиморфные параметры
    \end{itemize}


    Реализовано набором словарей, поддерживаемым при выводе типов:
    \begin{itemize}
        \item Объявленные типы (\texttt\small{list, map, option})
        \item Объявленные конструкторы (\texttt\small{Cons, Nil, Some, None})
        \item Схемы типов объявленных значений
    \end{itemize}
\end{frame}


\begin{frame}[fragile]
    \frametitle{Пример типизации ADT}
    \begin{lstlisting}[language=ML]
$ run << EOF
> type 'a option =
>   | Some of 'a
>   | None
>
> let rec find p ls =
>   match ls with
>   | [] -> None
>   | x :: xs -> if p x then Some x else find p xs

type '_0 option =
  | Some of '_0
  | None

let rec find: ('_4 -> bool) -> '_4 list -> '_4 option
		\end{lstlisting}
\end{frame}


\begin{frame}[fragile]
    \frametitle{Дальнейшная обработка match и конструкторов}

    \begin{itemize}
        \item Для каждого конструктора добавляется уникальный в рамках типа тег
        \item  \texttt\small{match-with} разворачивается в \texttt\small{if-then-else}
        \item При сопоставлении с образцом сначала сравнивается тег конструктора, затем его аргументы


    \end{itemize}
    \begin{lstlisting}[language=ML]
    $ run << EOF
    > let is_empty ls =
    >   match ls with
    >   | [] -> true
    >   | _ -> false

    let is_empty ls =
      let temp1 = ls in
      let temp2 = get_tag temp1 in
        (if (temp2 = 0)
          then true
          else false)
    \end{lstlisting}
\end{frame}

\begin{frame}
    \frametitle{Порождение кода}
    \begin{itemize}
        \item Значения алгебраических типов представляются как блоки переменного размера в куче

        \item Хранится заголовок с тегом и арностью, после него размещаются аргументы конструктора

        \item Аллокация, обращение к полям заголовка и аргументам конструктора реализованы на языке C
    \end{itemize}
\end{frame}


\begin{frame}[fragile]
    \frametitle{Пример, проходящий все стадии компиляции}

    \begin{lstlisting}[language=ML]
let rec map f ls =
  match ls with
  | [] -> []
  | x :: xs -> f x :: map f xs

let rec print_int_list ls =
  match ls with
    | [] -> 0
    | x :: xs ->
      let t = print x in
        print_int_list xs

let pow2 x = x * x

let main = print_int_list (map pow2 [ 1; 2; 3; 4; 5 ])
    \end{lstlisting}

\end{frame}


\begin{frame}[fragile]
    \frametitle{Тестирование}
    \begin{itemize}
        \item cram тесты dune
              \begin{itemize}
                  \item Для деревьев пишутся pretty-printer'ы
                  \item Модуль компилятора или сообщает об ошибке, или человеко-проверяемо печатает дерево
              \end{itemize}
        \item Процент покрытия измеряется с помощью утилиты \texttt\small{bisect\_ppx}
        \item В основном тестировались самописные реализации функций \texttt\small{Stdlib.List}
    \end{itemize}
\end{frame}


\begin{frame}
    \frametitle{Что ещё можно сделать}
    \begin{itemize}
        \item CPS convertion для \texttt\small{match-with} и конструкторов
        \item Преобразование \texttt\small{match-with} в деревья выбора или автоматы с возвратами
        \item Проверка полноты и неизбыточности паттерн-матчинга
    \end{itemize}
\end{frame}


\begin{frame}[fragile]
    \frametitle{Результаты}
    \begin{itemize}
        \item Окружение тайпчекера переработано для поддержки ADT
        \item Парсятся и тайпчекаются объявления типов, \texttt\small{match-with} и конструкторы ADT
        \item Добавлен сахар для конструкторов списков, их отдельная обработка в pretty-printer'ах и явное описание \texttt\small{TypeList} в \texttt\small{typedtree}
        \item Написан кодген \texttt\small{amd64} для значений алгебраических типов
    \end{itemize}
\end{frame}


\end{document}
